% ==============================================================
%  Chapter 5: AI for Data Scientists
% ==============================================================
\chapter{AI for Data Scientists}
\label{ch:datascience}

Data science involves a mix of programming, statistics, domain knowledge, and
communication --- all areas where AI adds value at different points.
The key distinction: AI can help you move faster through the mechanics of data work,
but it cannot substitute for the judgement required to ask the right question,
design a valid experiment, or interpret results in context.

\section{Exploratory Data Analysis}
\label{sec:eda}

EDA is where AI provides the most immediate time savings.
Generating pandas, Polars, or R code for standard exploration tasks --- distributions,
correlations, missing value analysis --- takes seconds.

\textbf{EDA prompt patterns:}
\begin{itemize}
  \item \textit{``Given this DataFrame schema [list columns and dtypes], write an EDA script that: plots distributions of all numeric columns, computes a correlation matrix, identifies columns with more than 5\% missing values, and flags potential outliers using IQR.''}
  \item \textit{``Explain what this distribution shape [describe or paste plot] suggests about the data generating process.''}
  \item \textit{``What are the standard statistical tests to check before assuming linearity between X and Y?''}
\end{itemize}

\begin{tipbox}
Use ChatGPT Code Interpreter or Julius by uploading your CSV directly.
These tools see your actual data and generate analysis code that runs immediately,
without you needing to describe the schema.
\end{tipbox}

\section{Feature Engineering}
\label{sec:features}

Feature engineering is part domain knowledge, part technique.
AI knows the technique well; you supply the domain knowledge.

\textbf{Effective feature engineering prompts:}
\begin{itemize}
  \item \textit{``I'm predicting customer churn for a SaaS product. My features include: login frequency, support ticket count, MRR, and contract length. What derived features and interactions would you add?''}
  \item \textit{``How should I encode this high-cardinality categorical feature (800 categories) for a gradient boosting model?''}
  \item \textit{``What are the standard methods for handling class imbalance in this binary classification problem (98\% negative class)?''}
\end{itemize}

AI generates the code for encoding, scaling, and feature interaction quickly.
You validate that the engineered features make sense in your business domain.

\section{Model Selection and Development}
\label{sec:models}

AI is excellent at explaining the tradeoffs between modelling approaches,
but it cannot run cross-validation on your data or know your latency/accuracy tradeoffs.

\textbf{Model selection prompts:}
\begin{itemize}
  \item \textit{``I have 50,000 labelled samples, 150 features (mixed numeric and categorical), a binary target, and a requirement for sub-50 ms inference. Compare XGBoost, LightGBM, and a logistic regression baseline.''}
  \item \textit{``What hyperparameters matter most for XGBoost on a high-cardinality, sparse dataset?''}
  \item \textit{``Write a scikit-learn pipeline that: encodes categoricals, scales numerics, trains an XGBoost classifier, and produces a calibrated probability output.''}
\end{itemize}

\begin{warnbox}
AI can generate training code but cannot validate that your train/test split is free of
leakage, that your evaluation metric matches the business objective, or that the
model will generalise to your production data distribution.
These require your domain knowledge.
\end{warnbox}

\section{Statistical Reasoning and Experiment Design}
\label{sec:stats}

AI is a surprisingly strong statistical reasoning partner.
It can explain test selection, identify violations of test assumptions, and structure
A/B test designs.

\textbf{Useful statistical prompts:}
\begin{itemize}
  \item \textit{``I want to test whether a new feature increased conversion rate. My current rate is 4.5\%, I want to detect a 0.5 pp improvement, with 80\% power and 5\% significance. What sample size do I need per group?''}
  \item \textit{``What assumptions does a two-sample t-test make, and how do I check each one in Python?''}
  \item \textit{``My data has heteroskedasticity. What regression alternatives should I consider?''}
\end{itemize}

\section{Interpreting and Communicating Results}
\label{sec:interpret}

Communicating findings to non-technical stakeholders is often harder than the analysis.
AI helps translate technical results into business language.

\textbf{Communication prompts:}
\begin{itemize}
  \item \textit{``I found a statistically significant increase in conversion rate (4.5\% to 5.1\%, p=0.003) from an A/B test. Write a two-paragraph summary for a non-technical VP.''}
  \item \textit{``Explain what a SHAP value is to a product manager with no ML background.''}
  \item \textit{``What are the limitations of this analysis I should proactively communicate?''}
\end{itemize}

\begin{tipbox}
Ask AI to steelman potential objections to your analysis.
\textit{``What would a critical statistician say is wrong with this methodology?''}
This surfaces weaknesses before stakeholders do.
\end{tipbox}
